% TOP_PROBLEM: {"id": 30007010, "title": "All rational-boundary projective lc fourfold programs terminate", "category_id": 5, "category": {"id": 5, "name": "algebraic_geometry", "display_name": "Algebraic Geometry", "description": "Geometric objects defined by polynomial equations.", "slug": "algebraic-geometry", "order_index": 5, "created_at": "2026-07-31T15:26:25.671Z"}, "status": "open"}
% FIRST_POSTED: 2026-10-06
% ATTEMPT_STATUS: solved
% OPENAI_CLAIM: {"schema_version":1,"source_repo":"https://github.com/openai/math","source_commit":"adc7f1241b42e322a6451854ab7e4b4c146bf78a","release_date":"2026-10-06","match":"full","resolution":"proved","independently_reviewed":false,"second_pass":"agreed","families":[{"family":"056","title":"Termination of projective and Kähler fourfold minimal model programs","subject":"Algebraic and complex geometry","manuscripts":[{"title":"Termination for projective log canonical fourfolds with rational boundary","pdf_path":"preprints/Termination-for-projective-log-canonical-fourfolds-with-rational-boundary-September-24-2026/paper.pdf","readme_path":"preprints/Termination-for-projective-log-canonical-fourfolds-with-rational-boundary-September-24-2026/README.md","date":"2026-09-24","theorem_ref":"thm:main (sections/00-introduction.tex:97)","statement_sources":[{"path":"preprints/Termination-for-projective-log-canonical-fourfolds-with-rational-boundary-September-24-2026/build/sections/00-introduction.tex","line":97,"label":"thm:main"}]}],"lean":null,"reasoning_trace":null}]}
\documentclass[11pt]{article}
\usepackage[margin=25mm]{geometry}
\usepackage{fontspec}
\setmainfont{Latin Modern Roman}
\usepackage{amsmath,amssymb,mathtools}
\usepackage{unicode-math}
\setmathfont{Latin Modern Math}
\usepackage{enumitem}
\usepackage{xurl,hyperref,bookmark}
\hypersetup{colorlinks=true,urlcolor=blue}
\setcounter{secnumdepth}{0}
\setlength{\parindent}{0pt}
\setlength{\parskip}{5pt}
\setlength{\emergencystretch}{3em}
\begin{document}
\section*{All rational-boundary projective lc fourfold programs terminate}

\subsection{OpenAI's claim}
Every permitted rational-boundary log-canonical fourfold program in characteristic zero is claimed to terminate, including mixed steps and arbitrary ray choices on initially nonfactorial models. The cited manuscript argues for that result; its relation to this problem is specified below.

\subsection{Mathematical statement of the claimed result}
\paragraph{Mathematical claim 1.} For a projective log canonical fourfold pair $(X,B)$ with rational boundary over a field that has characteristic zero and is algebraically closed, OpenAI claims termination of every permitted negative-extremal-ray program, without rational-factoriality or pseudo-effectivity assumptions. Its permitted birational steps are divisorial contractions, flips, or a divisorial contraction followed by a small positive modification; extra blowups and flops are excluded. It claims all required contractions and positive models exist at each stage, and every maximal program ends with
\[
K_Y+B_Y\text{ nef}\quad\text{or}\quad Y\to T\text{ a Mori fiber space}.
\]
The fiber contraction has connected fibers and normal base of smaller dimension, with its relative Picard number equal to one and its negative adjoint relatively ample. Nef inputs give a zero-step program.

\subsection{Outline of the argument}
Manuscript 1 (family 056, 2026-09-24): Divisor-cycle ranks discard all steps that lose divisors, and the remaining small tail lifts to crepant terminal junctions. A fixed set of boundary labels is constrained by surface discrepancy ACC; mixed-Hodge support bounds and finite model comparisons make the remaining numerical changes finite.

\subsection{Scope relative to this problem}
Manuscript 1 (family 056, 2026-09-24): The problem selects the mathematical statement summarized here, including its stated hypotheses and parameter range. The cited theorem gives a claimed proof of that assertion in full.

\subsection{Formal verification}
Family 056 has no Lean scope page in this snapshot; no formalization of these manuscript results is listed.

This site has not executed the Lean files or checked their dependencies. These descriptions report the scope of the linked declarations, rather than a new verification of the proofs.

\subsection{Status caveat}
OpenAI attributes these results to an internal model. Its release notice cautions that results lacking full formalization can contain mistakes. This site has not independently assessed the proofs; a solved label records the scope of the reported claim.

\subsection{Sources}
\begin{itemize}
\item OpenAI, \emph{Termination for projective log canonical fourfolds with rational boundary}, 2026-09-24; citation key \texttt{OAI:Termination-for-projective-log-canonical-fourfolds-with-rational-boundary-September-24-2026}. \href{https://github.com/openai/math/blob/adc7f1241b42e322a6451854ab7e4b4c146bf78a/preprints/Termination-for-projective-log-canonical-fourfolds-with-rational-boundary-September-24-2026/paper.pdf}{PDF (GitHub)}; \href{https://github.com/openai/math/raw/adc7f1241b42e322a6451854ab7e4b4c146bf78a/preprints/Termination-for-projective-log-canonical-fourfolds-with-rational-boundary-September-24-2026/paper.pdf}{PDF (direct)}; \href{https://github.com/openai/math/blob/adc7f1241b42e322a6451854ab7e4b4c146bf78a/preprints/Termination-for-projective-log-canonical-fourfolds-with-rational-boundary-September-24-2026/README.md}{Manuscript page and citation}; \href{https://github.com/openai/math/blob/main/preprints/Termination-for-projective-log-canonical-fourfolds-with-rational-boundary-September-24-2026/paper.pdf}{Latest version}; Statement source: \href{https://github.com/openai/math/blob/adc7f1241b42e322a6451854ab7e4b4c146bf78a/preprints/Termination-for-projective-log-canonical-fourfolds-with-rational-boundary-September-24-2026/build/sections/00-introduction.tex\#L97}{thm:main}.
\item \href{https://github.com/openai/math/blob/adc7f1241b42e322a6451854ab7e4b4c146bf78a/CONTENTS.md}{OpenAI catalogue entry}.
\end{itemize}
\end{document}
